\section{General Problem and Relevance}

Schon in den 1980er Jahren identifizierte die Forschung zu Group Support Systems (GSS) zentrale
Herausforderungen, die bis heute in unproduktiven Meetings wiederzufinden sind: ein unpassendes Tempo,
mangelhaftes Meeting-Design, fehlender Fokus, ausbleibende Entscheidungen oder ineffiziente Prozesse
\cite[p. 152]{beranekFacilitationGroupSupport1993}

Trotz des langjährigen Wissens um diese Probleme und der anerkannten Bedeutung von Facilitation-Ansätzen
scheinen diese Erkenntnisse in der Praxis noch immer nicht ausreichend umgesetzt zu werden.
Eine mögliche Erklärung liegt in der Kombination aus unzureichender Schulung, mangelnder Erfahrung und einer
natürlichen Resistenz gegenüber Veränderungen. \cite[p. 153]{hayneFacilitatorsPerspectiveMeetings}

Hinzu kommt, dass die Aufgaben eines Facilitators äußerst umfangreich und komplex sind und vermutlich nicht
jede Person für eine solche Rolle geeignet ist. [TODO: Quelle ergänzen]

KI-gestüzte Meeting Tools hätten das Potnezial diese Facilitating-Wissen demokratisieren.

Aktuelle Meeting-Tools unterstützen primär die Vorbereitung und Nachbereitung von Meetings. Funktionen
während des Meetings, die unter Facilitation kategorisiert werden könnten, sind dagegen eher reaktiv oder
interaktiv designt. Dies lässt sich unteranderem damit begründen, dass ein Großteil der Literatur darauf hin
deutet, dass KI-Systeme in diesem Kontext vor allem interaktiv gestaltet werden sollten und prokativer
ansätze deutlich schwieriger umzusetzen sind
(vgl.~\ref{sec:interaction-paradigms} und~\ref{sec:proactive-ai-interactions}).

Gleichzeitig treten in Meetings Probleme auf, die den Teilnehmenden häufig nicht bewusst sind, etwa
kognitive Verzerrungen, ein Abdriften vom Thema oder Kommunikationsschwierigkeiten. In solchen Fällen können
die Beteiligten ein KI-System nicht aktiv um Unterstützung bitten, weil sie sich der Problematik zu diesem
Zeitpunkt nicht bewusst sind. Auch wenn das System selbst einen Hinweis sendet, etwa über den Chat, kann
dieser ignoriert oder rationalisiert werden. \cite["4.4.3
Intervene"]{houttiObserveAskIntervene2025}

In der Human-Computer Interaction (HCI) existieren erste Arbeiten zu KI-Systemen in diesem Kontext. Der
überwiegende Teil bleibt jedoch theoretisch oder folgt einem interaktiven oder reaktiven Interaktionsmodell.
Praxisorientierte Systeme, die proaktiv und moderierend über Stimme in Meetings eingreifen, sind in der
Recherche kaum aufgefallen.

Diese Arbeit konzentriert sich daher auf konvergente Entscheidungsprozesse in Remote-Meetings, also auf
Gruppenprozesse, in denen vorhandene Informationen und Optionen bewertet und auf ein gemeinsames Ergebnis
verdichtet werden.
Diese Eingrenzung folgt dem Stand der bestehenden Meeting-Unterstützung: Vorbereitung und Nachbereitung sind in
kommerziellen Meeting-Tools bereits teilweise abgedeckt, und für divergente Praktiken (wie z.B. Brainstorming) liegen
vergleichsweise etablierte KI-Ansätze vor. [TODO:Quellen anführen]
In der konvergenten Phase bleiben zentrale GSS-Probleme wie fehlender Fokus, ausbleibender Abschluss und
verzerrter Informationsaustausch bestehen. Da sie den Teilnehmenden oft nicht bewusst sind, greift ein rein
interaktives Design hier zu kurz. Untersucht wird deshalb, wie ein KI-System in dieser Phase proaktiv moderieren
kann, um die Teilnehmenden zu unterstützen.

%%%%%%%%%%%%%%

% Die aktullen meetings-tools untersützten primär in der vorberitung und nachberitung des meetings udn Features
% die unter Facilitatrfunktion bezichnet werden könne ist eher reaktiv oder interactiv.

% viel der litarautr deutet darauf hin KI-System in diesem kontext eher interaktiv zu gestalten. Dies scheint
% im interesse der Meetingtiel nehmenden zu sein, weil diese der KI eine aktie Meeting untersützung nicht zutraune oder
% sorge haben in der hirachie übergangen zuwerden. Dem entsprechend gibt es in der Human-Computer-Interaction
% (HCI) erste ansätze die AI-Systeme in diesem kontext zu untersuchen. allerdings konnte weniger forschung
% dazugefunden werden die über theoretische Ansätze hinausgeht, wenn es um praxisorientierte KI-gestütze Meeting
% Support systeme geht, die proactiv in Meetings einschreiten.

% Bereits group support Systems (GSS) Forshcungs aus den 1980er Jahren hat einige probleme, welche zu
% unproductiven meetings führen identifiziert (z.B. Pace, Poor Meeting desing, Poor focus, Lack of closure,
% poor process usw.), die sich auch in modernen meetings immernoch wiederfinden.

% Trotz wissen über diese Probleme und der erkenniss über die wichtigkeit von facilitation ansätzen, scheinen
% dies erkenntnisse in der praxis nicht angekommen zu sein. Was vermutlich weiterhin auf unzreichendes trainung
% enerfahren heit und eine resisten gegen über veränderung zurückzuführen sein kann.
% \cite{beranekFacilitationGroupSupport1993}

% Reasearch bereits weitzurückrechende wissenscahftlich paper deuten darauf hin, das bereits minimaler
% facilitation support eine positiven effet auf die effectivität von gruppen meetings haben kann.
% \cite{hayneFacilitatorsPerspectiveMeetings}

% Current tools fall short in three key areas:
% \begin{itemize}
%   \item Passive functionality: Most AI tools (e.g., Zoom AI Companion, Google Meet’s recaps) focus on
%     post-meeting summaries or transcription, not real-time moderation [Zoom, 2024; Google, 2023].
%   \item Generic interventions: Solutions like Miro’s AI facilitator provide template-based guidance but fail
%     to adapt to the meeting’s context or user preferences [Miro, 2024].
%   \item Lack of empirical validation: Few studies evaluate user acceptance of proactive AI moderation,
%     particularly in high-stakes decision-making scenarios [Benbasat and Barki, 2007].
% \end{itemize}

% To address these gaps, this thesis proposes a proactive AI moderation system that:
% \begin{itemize}
%   \item Detects thematic deviations in real time using NLP and meeting goal tracking.
%   \item Intervenes constructively with personalized, non-intrusive prompts (e.g., private messages to
%     individuals or group reminders).
%   \item Validates design choices through user studies focusing on usability (SUS) and acceptance (TAM).
% \end{itemize}
% Such a system could reduce meeting fatigue, improve decision quality, and lower the barrier to effective
% facilitation—but its success hinges on user trust and perceived usefulness, which remain underexplored in literature.
